\chap{Field Effect Transistor}

\section{Introduction}
In this laboratory you study field-effect transistors: the characteristics of the JFET and of the depletion and enhancement MOSFETs, their DC bias circuits, and two applications of the MOSFET as a switch (a lamp and a DC motor).

\textbf{After this laboratory you will be able to:}
\begin{itemize}
  \item plot the output and transfer characteristics of a JFET and a MOSFET with DC sweeps and nested sweeps;
  \item calculate and verify the bias of self-bias and voltage-divider circuits;
  \item identify the triode and saturation regions from the simulated voltages;
  \item use a MOSFET as a switch and explain the role of the flyback diode;
  \item measure averages and ripples of transient waveforms with \cmd{.meas}.
\end{itemize}

\textbf{Exercises.} Save every schematic as \file{Lab5/ExNN.asc} (\file{Ex01.asc} to \file{Ex10.asc} for main exercises; \file{Ex11.asc} and \file{Ex12.asc} for advanced exercises). Calculation spaces are blank; fill in the comparison tables with your own theoretical and simulated values.

\section{LTspice guide}

\subsection{JFET and MOSFET symbols}
Press \texttt{F2} and search for \cmd{njf} (n-channel JFET), \cmd{pjf}, \cmd{nmos} (n-channel MOSFET) or \cmd{pmos}. The MOSFET symbol has four terminals (drain, gate, source and bulk): connect the bulk to the source. Declare the model with a SPICE directive and type its name in the value field of the transistor; for a MOSFET the instance parameters \cmd{L} and \cmd{W} are entered in the \textit{SpiceLine} (right-click the symbol), and they are set to equal values here so that $K_p$ gives the transconductance directly.

\labfig{lab5_guide01.png}{10cm}{Placing a JFET and declaring its model}{fig:lab5_guide01}

\subsection{Nested DC sweeps for FET characteristics}
The family of output characteristics is plotted with a nested DC sweep of the drain-source voltage and of the gate-source voltage (see the LTspice guide of Lab~3). Plot the drain current by clicking on the transistor symbol in the schematic.

\labfig{lab5_guide02.png}{12cm}{Nested DC sweep for a JFET}{fig:lab5_guide02}

\subsection{Pulse sources and PWM}
A PWM gate drive is a pulse source: in the voltage source dialog click \textit{Advanced} and choose PULSE, with the low value, the high value, the delay, the rise and fall times, the on time and the period. The average of a signal is measured with \cmd{.meas TRAN name AVG expression FROM=t1 TO=t2}.

\section{Exercises}

\subsection{Exercise 1: JFET output characteristic}
\status{Unchanged (Ch.~6: JFET output characteristic with $V_{GS}=0$).}
In the circuit of Figure~\ref{fig:lab5_ex01} the gate of an n-channel JFET is connected to the source ($V_{GS}=0$) and the drain voltage $V_{DS}$ is swept. The JFET has $I_{DSS}=10$\,mA and $V_P=-2$\,V ($I_D=I_{DSS}(1-V_{GS}/V_P)^2$ in saturation). Sketch the expected curve $I_D=f(V_{DS})$, indicate the ohmic region and the saturation region and the value of $V_{DS}$ at the boundary. Then plot the curve with a DC sweep of $V_{DS}$ from 0 to 10\,V and read $I_{DSS}$.

\labfig{lab5_ex01.png}{8cm}{Exercise 1: JFET with $V_{GS}=0$}{fig:lab5_ex01}

\theory{6}

\simulation{
\item Draw the circuit of Figure~\ref{fig:lab5_ex01} in LTspice and save it as \file{Lab5/Ex01.asc}.
\item Declare the JFET model with a directive (\texttt{S} key): \cmd{.model JN NJF(Beta=2.5m Vto=-2 Lambda=0.01)} ($I_{DSS}=\beta V_P^2=10$\,mA, $V_P=-2$\,V). Place the transistor (\cmd{njf}) and type \texttt{JN} in its value field.
\item Run \cmd{.dc VDS 0 10 0.05} and plot \cmd{Id(J1)} (click the JFET).
\item Read the saturation current and the knee voltage with the cursors.
}

\comparetable{
\cmprow{$I_{DSS}$ (mA)}
\cmprow{Boundary $V_{DS}$ (V)}
\cmprow{Slope in the ohmic region (mS)}
}

\commentlines{2}

\subsection{Exercise 2: JFET family of curves (nested DC sweep)}
\status{Unchanged (Ch.~6: JFET family of curves using nested DC sweep).}
Use the circuit of Figure~\ref{fig:lab5_ex02} to plot the family of output characteristics of the JFET for $V_{GS}=-2$, $-1.5$, $-1$, $-0.5$ and 0\,V with a nested DC sweep. The JFET has $I_{DSS}=10$\,mA and $V_P=-2$\,V ($I_D=I_{DSS}(1-V_{GS}/V_P)^2$ in saturation). Calculate the saturation current for each $V_{GS}$ and compare it with the curves.

\labfig{lab5_ex02.png}{8cm}{Exercise 2: JFET with variable $V_{GS}$ and $V_{DS}$}{fig:lab5_ex02}

\theory{8}

\simulation{
\item Draw the circuit of Figure~\ref{fig:lab5_ex02} in LTspice and save it as \file{Lab5/Ex02.asc}.
\item Declare the JFET model with a directive (\texttt{S} key): \cmd{.model JN NJF(Beta=2.5m Vto=-2 Lambda=0.01)} ($I_{DSS}=\beta V_P^2=10$\,mA, $V_P=-2$\,V). Place the transistor (\cmd{njf}) and type \texttt{JN} in its value field.
\item Use the nested sweep \cmd{.dc VDS 0 10 0.05 VGS -2 0 0.5} (first source $V_{DS}$, second source $V_{GS}$) and plot \cmd{Id(J1)}.
\item Read the saturation current of each curve at $V_{DS}=8$\,V.
}

\comparetable{
\cmprow{$I_D$ at $V_{GS}=0$ (mA)}
\cmprow{$I_D$ at $V_{GS}=-0.5$\,V (mA)}
\cmprow{$I_D$ at $V_{GS}=-1$\,V (mA)}
\cmprow{$I_D$ at $V_{GS}=-1.5$\,V (mA)}
\cmprow{$I_D$ at $V_{GS}=-2$\,V (mA)}
}

\commentlines{2}

\subsection{Exercise 3: JFET transfer characteristic}
\status{Supplementary (JFET transfer characteristic and transconductance $g_m$).}
In the circuit of Figure~\ref{fig:lab5_ex03} the drain voltage is fixed at 8\,V and the gate-source voltage is swept. The JFET has $I_{DSS}=10$\,mA and $V_P=-2$\,V ($I_D=I_{DSS}(1-V_{GS}/V_P)^2$ in saturation). Calculate the transconductance $g_m=\partial I_D/\partial V_{GS}$ at $V_{GS}=-1$\,V and at $V_{GS}=-0.5$\,V, then plot the transfer characteristic $I_D=f(V_{GS})$ and the pinch-off voltage.

\labfig{lab5_ex03.png}{8cm}{Exercise 3: JFET transfer characteristic}{fig:lab5_ex03}

\theory{8}

\simulation{
\item Draw the circuit of Figure~\ref{fig:lab5_ex03} in LTspice and save it as \file{Lab5/Ex03.asc}.
\item Declare the JFET model with a directive (\texttt{S} key): \cmd{.model JN NJF(Beta=2.5m Vto=-2 Lambda=0.01)} ($I_{DSS}=\beta V_P^2=10$\,mA, $V_P=-2$\,V). Place the transistor (\cmd{njf}) and type \texttt{JN} in its value field.
\item Run \cmd{.dc VGS -3 0 0.01} and plot \cmd{Id(J1)}.
\item Plot the derivative \cmd{d(Id(J1))} (right-click the plot window $\rightarrow$ \textit{Add Trace}, expression \cmd{d(Id(J1))}) or estimate $g_m$ from two points.
}

\comparetable{
\cmprow{$V_P$ (V)}
\cmprow{$I_D$ at $V_{GS}=-1$\,V (mA)}
\cmprow{$g_m$ at $V_{GS}=-1$\,V (mS)}
\cmprow{$g_m$ at $V_{GS}=-0.5$\,V (mS)}
}

\commentlines{2}

\subsection{Exercise 4: JFET self-bias}
\status{Modified (Ch.~6: JFET self-bias circuit analysis).}
In the self-bias circuit of Figure~\ref{fig:lab5_ex04} (V$_{DD}$ = 15\,V, R$_D$ = 1.2\,k$\Omega$, R$_G$ = 1\,M$\Omega$, R$_S$ = 470\,$\Omega$) determine the quiescent values $I_D$, $V_{GS}$, $V_{DS}$ and $V_D$. The JFET has $I_{DSS}=10$\,mA and $V_P=-2$\,V ($I_D=I_{DSS}(1-V_{GS}/V_P)^2$ in saturation). (Write $V_{GS}=-I_DR_S$ and solve the resulting quadratic equation for $I_D$.) Check that the JFET works in saturation. Then simulate.

\labfig{lab5_ex04.png}{8cm}{Exercise 4: self-biased JFET}{fig:lab5_ex04}

\theory{10}

\simulation{
\item Draw the circuit of Figure~\ref{fig:lab5_ex04} in LTspice and save it as \file{Lab5/Ex04.asc}.
\item Declare the JFET model with a directive (\texttt{S} key): \cmd{.model JN NJF(Beta=2.5m Vto=-2 Lambda=0.01)} ($I_{DSS}=\beta V_P^2=10$\,mA, $V_P=-2$\,V). Place the transistor (\cmd{njf}) and type \texttt{JN} in its value field.
\item Add the \cmd{.op} analysis (\textit{Simulate $\rightarrow$ Edit Simulation Cmd $\rightarrow$ DC op pnt}), run the simulation and read the node voltages and branch currents.
\item Be careful with signs: LTspice reports the current of a voltage source as positive when it flows \emph{into} its positive terminal, and the current of a two-terminal device as positive when it flows from its first pin to its second pin.
}

\comparetable{
\cmprow{$I_D$ (mA)}
\cmprow{$V_{GS}$ (V)}
\cmprow{$V_D$ (V)}
\cmprow{$V_{DS}$ (V)}
}

\commentlines{2}

\subsection{Exercise 5: JFET voltage-divider bias}
\status{Modified (Ch.~6: JFET voltage-divider bias circuit analysis).}
In the circuit of Figure~\ref{fig:lab5_ex05} (V$_{DD}$ = 15\,V, R1 = 1\,M$\Omega$, R2 = 270\,k$\Omega$, R$_D$ = 2.2\,k$\Omega$, R$_S$ = 1.5\,k$\Omega$) determine the gate voltage $V_G$, and the quiescent values $I_D$, $V_{GS}$, $V_S$, $V_D$ and $V_{DS}$. The JFET has $I_{DSS}=10$\,mA and $V_P=-2$\,V ($I_D=I_{DSS}(1-V_{GS}/V_P)^2$ in saturation). Compare this bias circuit with the self-bias circuit of Exercise~4 and explain why it makes $I_D$ less dependent on the JFET parameters. Verify by simulation.

\labfig{lab5_ex05.png}{8cm}{Exercise 5: JFET with a voltage divider}{fig:lab5_ex05}

\theory{10}

\simulation{
\item Draw the circuit of Figure~\ref{fig:lab5_ex05} in LTspice and save it as \file{Lab5/Ex05.asc}.
\item Declare the JFET model with a directive (\texttt{S} key): \cmd{.model JN NJF(Beta=2.5m Vto=-2 Lambda=0.01)} ($I_{DSS}=\beta V_P^2=10$\,mA, $V_P=-2$\,V). Place the transistor (\cmd{njf}) and type \texttt{JN} in its value field.
\item Add the \cmd{.op} analysis (\textit{Simulate $\rightarrow$ Edit Simulation Cmd $\rightarrow$ DC op pnt}), run the simulation and read the node voltages and branch currents.
\item Be careful with signs: LTspice reports the current of a voltage source as positive when it flows \emph{into} its positive terminal, and the current of a two-terminal device as positive when it flows from its first pin to its second pin.
}

\comparetable{
\cmprow{$V_G$ (V)}
\cmprow{$I_D$ (mA)}
\cmprow{$V_{GS}$ (V)}
\cmprow{$V_S$ (V)}
\cmprow{$V_D$ (V)}
\cmprow{$V_{DS}$ (V)}
}

\commentlines{2}

\subsection{Exercise 6: Depletion-mode MOSFET}
\status{Modified (Ch.~6: depletion-mode MOSFET self-bias and model declaration).}
A depletion-mode n-channel MOSFET has $I_{DSS}=8$\,mA and a threshold voltage $V_T=-3$\,V ($I_D=I_{DSS}(1-V_{GS}/V_T)^2$). In the self-bias circuit of Figure~\ref{fig:lab5_ex06} (V$_{DD}$ = 12\,V, R$_D$ = 1\,k$\Omega$, R$_G$ = 1\,M$\Omega$, R$_S$ = 500\,$\Omega$) determine $I_D$, $V_{GS}$ and $V_{DS}$ and check the saturation condition. Then simulate. Explain why a depletion MOSFET can work with a negative or a positive $V_{GS}$.

\labfig{lab5_ex06.png}{8cm}{Exercise 6: self-biased depletion MOSFET}{fig:lab5_ex06}

\theory{10}

\simulation{
\item Draw the circuit of Figure~\ref{fig:lab5_ex06} in LTspice and save it as \file{Lab5/Ex06.asc}.
\item Declare the model with a negative threshold: \cmd{.model DMOS NMOS(Kp=1.78m Vto=-3 Lambda=0.01)} ($K_p=2I_{DSS}/V_T^2$). Place \cmd{nmos}, connect its bulk to the source and set the instance parameters to $W=L$ (for example \cmd{L=1u W=1u}), so that $K_p$ is used directly.
\item Add the \cmd{.op} analysis (\textit{Simulate $\rightarrow$ Edit Simulation Cmd $\rightarrow$ DC op pnt}), run the simulation and read the node voltages and branch currents.
\item Be careful with signs: LTspice reports the current of a voltage source as positive when it flows \emph{into} its positive terminal, and the current of a two-terminal device as positive when it flows from its first pin to its second pin.
}

\comparetable{
\cmprow{$I_D$ (mA)}
\cmprow{$V_{GS}$ (V)}
\cmprow{$V_D$ (V)}
\cmprow{$V_{DS}$ (V)}
}

\commentlines{2}

\subsection{Exercise 7: Enhancement MOSFET characteristics}
\status{Modified (Ch.~6: enhancement-mode MOSFET output characteristics).}
An n-channel enhancement MOSFET has a threshold voltage $V_T=2$\,V and $I_D=K(V_{GS}-V_T)^2$ in saturation with $K=1$\,mA/V$^2$ (and $I_D=2K[(V_{GS}-V_T)V_{DS}-V_{DS}^2/2]$ in the triode region). Using Figure~\ref{fig:lab5_ex07}, plot the output characteristics for $V_{GS}=3$, 4, 5 and 6\,V with a nested DC sweep. Calculate the saturation current for each $V_{GS}$ and the boundary value of $V_{DS}$ between the triode and the saturation region, and compare with the curves.

\labfig{lab5_ex07.png}{8cm}{Exercise 7: enhancement MOSFET characteristics}{fig:lab5_ex07}

\theory{8}

\simulation{
\item Draw the circuit of Figure~\ref{fig:lab5_ex07} in LTspice and save it as \file{Lab5/Ex07.asc}.
\item Declare \cmd{.model NM NMOS(Kp=2m Vto=2 Lambda=0.01)} and set the instance parameters to $W=L$ (for example \cmd{L=1u W=1u}).
\item Run the nested sweep \cmd{.dc VDS 0 10 0.05 VGS 3 6 1} and plot \cmd{Id(M1)}.
\item Read the saturation current of each curve at $V_{DS}=8$\,V.
}

\comparetable{
\cmprow{$I_D$, $V_{GS}=3$\,V (mA)}
\cmprow{$I_D$, $V_{GS}=4$\,V (mA)}
\cmprow{$I_D$, $V_{GS}=5$\,V (mA)}
\cmprow{$I_D$, $V_{GS}=6$\,V (mA)}
}

\commentlines{2}

\subsection{Exercise 8: Enhancement MOSFET bias}
\status{Supplementary (enhancement-mode MOSFET voltage-divider bias circuit).}
In the voltage-divider bias circuit of Figure~\ref{fig:lab5_ex08} (V$_{DD}$ = 12\,V, R1 = 1.5\,M$\Omega$, R2 = 1\,M$\Omega$, R$_D$ = 2.2\,k$\Omega$, R$_S$ = 1\,k$\Omega$; $V_T=2$\,V, $K=1$\,mA/V$^2$) determine $V_G$, $I_D$, $V_{GS}$, $V_D$ and $V_{DS}$ and check the saturation condition $V_{DS}>V_{GS}-V_T$. Then simulate.

\labfig{lab5_ex08.png}{8cm}{Exercise 8: enhancement MOSFET with a voltage divider}{fig:lab5_ex08}

\theory{10}

\simulation{
\item Draw the circuit of Figure~\ref{fig:lab5_ex08} in LTspice and save it as \file{Lab5/Ex08.asc}.
\item Declare \cmd{.model NM NMOS(Kp=2m Vto=2 Lambda=0.01)} and set \cmd{L=1u W=1u} in the instance parameters.
\item Add the \cmd{.op} analysis (\textit{Simulate $\rightarrow$ Edit Simulation Cmd $\rightarrow$ DC op pnt}), run the simulation and read the node voltages and branch currents.
\item Be careful with signs: LTspice reports the current of a voltage source as positive when it flows \emph{into} its positive terminal, and the current of a two-terminal device as positive when it flows from its first pin to its second pin.
}

\comparetable{
\cmprow{$V_G$ (V)}
\cmprow{$I_D$ (mA)}
\cmprow{$V_{GS}$ (V)}
\cmprow{$V_S$ (V)}
\cmprow{$V_D$ (V)}
\cmprow{$V_{DS}$ (V)}
}

\commentlines{2}

\subsection{Exercise 9: MOSFET as a switch}
\status{Modified (Ch.~6: MOSFET as a power switch).}
A MOSFET can switch a load such as a lamp (modelled by a 47\,$\Omega$ resistor connected to a 12\,V supply, see Figure~\ref{fig:lab5_ex09}). The gate is driven by a 0\,V/5\,V pulse train (period 20\,ms, 50\,\% duty cycle). The transistor has $V_T=2$\,V and $K_p=0.5$\,A/V$^2$ (on-resistance in the triode region $R_{DS(on)}\approx1/[K_p(V_{GS}-V_T)]$). Calculate $R_{DS(on)}$, the load current, the drain voltage and the power dissipated in the transistor when it is on, and the drain voltage when it is off. Then run a transient simulation of 60\,ms.

\labfig{lab5_ex09.png}{8cm}{Exercise 9: MOSFET switch}{fig:lab5_ex09}

\theory{10}

\simulation{
\item Draw the circuit of Figure~\ref{fig:lab5_ex09} in LTspice and save it as \file{Lab5/Ex09.asc}.
\item Declare \cmd{.model NMSW NMOS(Kp=0.5 Vto=2)} and set \cmd{L=1u W=1u}.
\item Set V2 to a pulse source (PULSE: 0\,V, 5\,V, 0 delay, 1\,$\micro$s rise and fall, 10\,ms on time, 20\,ms period) and run \cmd{.tran 0 60m 0 10u}.
\item Plot the gate voltage and the drain voltage and read the on-state drain voltage and the current in R$_L$.
}

\comparetable{
\cmprow{$R_{DS(on)}$ ($\Omega$)}
\cmprow{Load current when on (mA)}
\cmprow{Drain voltage when on (mV)}
\cmprow{Power in the MOSFET when on (mW)}
\cmprow{Drain voltage when off (V)}
}

\commentlines{2}

\subsection{Exercise 10: DC motor control with a MOSFET}
\status{Modified (Ch.~6: DC motor drive with flyback protection diode).}
A small DC motor is modelled in Figure~\ref{fig:lab5_ex10} by a 10\,$\Omega$ resistor in series with a 10\,mH inductor, supplied by 12\,V and switched by a MOSFET driven with a 1\,kHz PWM signal with a 50\,\% duty cycle. A diode D1 (flyback diode) is connected across the motor. Calculate the average voltage applied to the motor, the average motor current, and explain the role of the diode (what would happen to the drain voltage at turn-off without it?). Then run a transient simulation of 30\,ms.

\labfig{lab5_ex10.png}{10cm}{Exercise 10: motor drive with a flyback diode}{fig:lab5_ex10}

\theory{10}

\simulation{
\item Draw the circuit of Figure~\ref{fig:lab5_ex10} in LTspice and save it as \file{Lab5/Ex10.asc}.
\item Reuse the MOSFET model of Exercise~9. Drive the gate with PULSE (0\,V, 5\,V, 0 delay, 1\,$\micro$s rise and fall, 0.5\,ms on time, 1\,ms period). Declare a diode model \cmd{.model D1N4148 D(Is=2.52n Rs=0.568 N=1.752)}.
\item Run \cmd{.tran 0 30m 0 2u}. Plot the motor current \cmd{I(Lm)} and the drain voltage.
\item Measure the average current in steady state with \cmd{.meas TRAN Iavg AVG I(Lm) FROM=20m TO=30m}.
\item Remove the diode and run again: observe the drain voltage spikes.
}

\comparetable{
\cmprow{Average motor voltage (V)}
\cmprow{Average motor current (mA)}
\cmprow{Current ripple (mA)}
\cmprow{Peak drain voltage without diode (V)}
}

\commentlines{2}

\newpage
\section{Advanced exercises}

\subsection{Advanced Exercise 1: On-resistance of a MOSFET versus the gate voltage}
\status{Advanced / Supplementary (MOSFET on-resistance $R_{DS(on)}$ versus gate-source voltage).}
In the circuit of Figure~\ref{fig:lab5_adv01} the drain is held at a small voltage (0.1\,V) so that the MOSFET works in the triode region. Study the on-resistance $R_{DS(on)}=V_{DS}/I_D$ as a function of the gate-source voltage.

\labfig{lab5_adv01.png}{8cm}{Advanced Exercise 1: MOSFET in the triode region}{fig:lab5_adv01}

\subsubsection{Theory calculation}\vspace{-3pt}\notesblock Using $R_{DS(on)}\approx1/[K_p(V_{GS}-V_T)]$ with $K_p=0.5$\,A/V$^2$ and $V_T=2$\,V, calculate $R_{DS(on)}$ for $V_{GS}=3$, 5 and 10\,V.\dotlines{4}

\simulation{
\item Draw the circuit of Figure~\ref{fig:lab5_adv01} in LTspice and save it as \file{Lab5/Ex11.asc}.
\item Run \cmd{.dc VGS 0 10 0.05} and plot \cmd{0.1/Id(M1)} (right-click the plot $\rightarrow$ \textit{Add Trace}). Limit the vertical axis to see the useful range.
\item Read $R_{DS(on)}$ at $V_{GS}=3$, 5 and 10\,V with the cursors.
}
\begin{center}\renewcommand{\arraystretch}{1.8}\begin{tabular}{|c|c|c|c|}\hline \textbf{$V_{GS}$ (V)} & \textbf{$R_{DS(on)}$ theory ($\Omega$)} & \textbf{$R_{DS(on)}$ LTspice ($\Omega$)} & \textbf{Error (\%)}\\ \hline 3 & & & \\ \hline 5 & & & \\ \hline 10 & & & \\ \hline \end{tabular}\end{center}
\commentlines{2}
\newpage

\subsection{Advanced Exercise 2: CMOS inverter}
\status{Advanced / Supplementary (CMOS inverter voltage transfer characteristic and switching threshold).}
Figure~\ref{fig:lab5_adv02} shows a CMOS inverter made of an n-channel MOSFET ($V_T=2$\,V, $K_p=2$\,mA/V$^2$) and a p-channel MOSFET ($V_T=-2$\,V, $K_p=1$\,mA/V$^2$), supplied by 5\,V. Plot its voltage transfer characteristic and determine the logic switching threshold.

\labfig{lab5_adv02.png}{8cm}{Advanced Exercise 2: CMOS inverter}{fig:lab5_adv02}

\subsubsection{Theory calculation}\vspace{-3pt}\notesblock Explain qualitatively which transistor conducts for $V_{in}=0$\,V and $V_{in}=5$\,V and estimate the input voltage at which both transistors are in saturation.\dotlines{4}

\simulation{
\item Draw the circuit of Figure~\ref{fig:lab5_adv02} in LTspice and save it as \file{Lab5/Ex12.asc}.
\item Declare \cmd{.model NM NMOS(Kp=2m Vto=2 Lambda=0.01)} and \cmd{.model PM PMOS(Kp=1m Vto=-2 Lambda=0.01)}; use \cmd{nmos} and \cmd{pmos} symbols with $W=L$.
\item Run \cmd{.dc Vin 0 5 0.01} and plot \cmd{V(out)}; read the threshold where $V_{out}=V_{in}$.
\item Plot the supply current \cmd{I(V1)} to see the current peak during the transition.
}
\begin{center}\renewcommand{\arraystretch}{1.8}\begin{tabular}{|c|c|c|c|}\hline \textbf{Quantity} & \textbf{Theory} & \textbf{LTspice} & \textbf{Error (\%)}\\ \hline Switching threshold (V) & & & \\ \hline Output for $V_{in}=0$\,V (V) & & & \\ \hline Output for $V_{in}=5$\,V (V) & & & \\ \hline \end{tabular}\end{center}
\commentlines{2}
\newpage
